\section{Hercules’ Sixth Task: The Stymphalian Birds}

Hercules now advances\footnote{\url{https://en.wikipedia.org/wiki/Stymphalian\_Birds}} to the sphere of war. Openly, now, following the altruistic stooping to help the needy and the ``folk'', Hercules comes bearing arms against the wicked. It is almost as if Hercules has doubly earned the right to tilt in combat, through having served as a squire to the filthy stables.

Mastering physical nature, he now is capable of turning his hands to subtle enemies, the enemies and spirits of the aether. He is being annointed to wage war against the very worst and most violent of such enemies. Having conquered the ``birds'' inside of him through humility and ingenuity and sheer sense, they are now available ``outside'' of him, clearly differentiated, as ``enemies''. The earlier labors were against the ``adversary'', whose job is to be the primary enemy and ``train'' the hero. The lion and the hydra were more primal threats to Hercules, roaring monsters who reflected the devouring urges of the Ego, expressed through sex and innumerable physical wants. These birds are exteriorized as ``lesser'' demons, but have less obvious uses for the hero. They are more like parasites, and so are more pestilential. They have no ``sacred'' use. They are simply a scourge, by their sheer numbers and vicious actions. If the lion and the hydra represented focused spiritual dangers close to the heart of the hero, these birds obviously incarnate the ``\emph{asuras}'' --- the ``birds of the air'' that pick up seeds that might otherwise sprout and flower in the hearts of the people. They are what most people think of when they speak the word ``demon''. They are on no account easily confused with the ``daemon'': by contrast, a lion or hydra or boar might be mistaken, in a state of spiritual confusion, for one's own daemon, or guiding angel.

Hercules brings war against the \emph{asuras} --- the man of Tradition does not mistake Lucifer for Satan, or Satan for the lesser imps. He sees chaos, and recognizes it for what it is. Likewise, he is clear-sighted enough to see that these creatures will not provide him with anything other than war and sport. He is the falcon, they are the quarry. \emph{There will be no ``cloak'' or ``blood'' that he obtains from contact with these beings, who are simply negative energy that needs to be banished from the ``center'' (I say the center, because all creatures have a ``place'')}. A great part of the confusion, even in the ranks of Tradition, stems from being unable to differentiate a bird from a lion. Indulging in drugs, practicing sex magic with random partners, and vociferously attacking the outer world indiscriminately with Leftist rage is a state that results from not having commenced in the proper order, and subdued the ``inner lion'' and the many-headed snake inside of one's self, first.

It is the chronic condition of the basic, inner, true Leftist. They are simply incapable of recognizing their own faults and basic condition, often on any level. When they do attempt this, they reduce every human being (presumably excluding themselves) to a cynical level of animal drives and primal urges. This is because they experience their own Self at this level, and so are incapable of imagining the existence of anything else in another being. This is the ``lowest common denominator'', by which they shift responsibility for their faults on to some kind of Deistic clock-work Creator, who made everyone equally worthless and mechanical/animal. The New Right runs afoul of this when they think that making everything more ``white'' would protect the West. It would, in an imaginary Leftist universe; unfortunately, the ``lion'' in the camp is actually a white lion, since the spearhead of anti-Tradition has a white face. Far better to work for a restoration of noble bloodlines and families, regardless of their ``color'', regardless of political ``status''. The longer this is put off, the less white such an inevitable Restoration would be. The answer is to stop thinking about skin color \emph{per se}, and to seek spiritual excellence in conformity with antique European Tradition, whose roots lie in Classicism, Egypt, and older worlds still. This would be the political equivalent of the First Labor.

By the Sixth Labor, we see the emergence of a new balance\footnote{\url{https://gornahoor.net/?p=5808}}: Life, properly itself, and manifest properly in Justice, emerges, as a perfect blending out of a perfect blending (2+3 = 5, but 2×3 = 6).. There is, in this Labor, as in the Hexad, something totally perfect and complete. The Sixth Labor shows the hero as he finally will be --- easily defeating dreadful adversaries, who are really no more than really obnoxious and odious predators, swiftly falling and fleeing before the ``Strong Man''. In the labor of the Pentad, Hercules attains perfection by uniting the two numbers that ``are not 1'', through addition: this leads to ``Life'' -- the folk hero willing to clean the stables, who has ample strength and pluck to easily fix a huge, stinking mess. In the labor of Hexad, this is multiplied and amplified into the spiritual sphere --- because he is ``Robin Hood'', therefore, he can also be the ``one true Englishman'' to the crown. He can defeat the imps.

The Sixth Labor is associated with Sagittarius in the Zodiac. After conquering Sin in Scorpio, the hero must master the power of right action and thought. He passes from triumph over the serpent, to mastery of insight. Thus, Hercules sees an easy and swift way to achieve victory over the birds, rather than killing each one laboriously. He sounds the cymbals, and drives them in droves and flocks to the air, killing many directly or indirectly. The rest flee.

\begin{quotex}
These fly against those who come to hunt them, wounding and killing them with their beaks. All armor of bronze or iron that men wear is pierced by the birds; but if they weave a garment of thick cork, the beaks of the Stymphalian birds are caught in the cork garment... These birds are of the size of a crane, and are like the ibis, but their beaks are more powerful, and not crooked like that of the ibis.
\flright{\textsc{Pausanias}, \emph{Description of Greece}, 8.22.5}
\end{quotex}


In Scripture, wild birds living in desolate places symbolize the \emph{asuras} -- the lesser demons and imps. Hercules is coming into his own, like an invincible hero who runs the course of the Sun. After defeating internal enemies, he is fit to ``manifest'' in the physical world as Life/Justice, then Perfection/Beauty, in this labor and the last.

Hercules, like the falcon of the tower or the hawk of the fist, has been trained in partnership as an ``equal'' to the gods, and is now set loose upon his quarry\footnote{\url{https://www.youtube.com/watch?v=Hdo0VnsJSxo}}. His instincts are purified, and he now wages war.


\flrightit{Posted on 2013-09-27 by Logres}

